\chapter{Additional Results}
\label{app:additional-results}

These results extend the comparisons in Chapter~\ref{chap:results}.
Each setting uses 20,000 paired samples at the two-sided 5\% threshold.

\section{Individual main-grid null results}

Table~\ref{tab:all-main-null} gives every main-grid null setting separately,
across four table sizes, three margin profiles and nine sample sizes.
All have \(I(P)=I(Q)=0.02\) nats. Uniform and same-skew settings have
\(P=Q\); different-skew settings have \(P\ne Q\) with equal MI\@.
The margins are specified in Table~\ref{tab:margin-profiles}.
Rejection rates divide valid rejections by all 20,000 samples; the valid-result
rate shows the fraction returning a usable p-value. At \(n=2\), both methods
have zero validity. Their zero rejection rates therefore indicate failed
calculations, not well-calibrated tests.

\begingroup\small
\setlength{\tabcolsep}{4pt}
\begin{longtable}{@{}llrrrrr@{}}
\caption{False-positive and valid-result rates for all 108 main null settings.}
\label{tab:all-main-null}\\
\toprule
Shape & Margins & $n$ & \multicolumn{2}{c}{Normal Wald} & \multicolumn{2}{c}{Expanded Welch}\\
\cmidrule(lr){4-5}\cmidrule(l){6-7}
 & & & Reject & Valid & Reject & Valid\\
\midrule\endfirsthead
\toprule
Shape & Margins & $n$ & \multicolumn{2}{c}{Normal Wald} & \multicolumn{2}{c}{Expanded Welch}\\
\cmidrule(lr){4-5}\cmidrule(l){6-7}
 & & & Reject & Valid & Reject & Valid\\
\midrule\endhead
\(2\times2\) & Uniform & 2 & 0.0000 & 0.0000 & 0.0000 & 0.0000 \\
\(2\times2\) & Uniform & 5 & 0.0914 & 0.9862 & 0.0000 & 0.7774 \\
\(2\times2\) & Uniform & 10 & 0.0588 & 0.9889 & 0.0263 & 0.9795 \\
\(2\times2\) & Uniform & 20 & 0.0261 & 0.9977 & 0.0164 & 0.9833 \\
\(2\times2\) & Uniform & 50 & 0.0082 & 0.9999 & 0.0045 & 0.9938 \\
\(2\times2\) & Uniform & 100 & 0.0142 & 1.0000 & 0.0063 & 0.9991 \\
\(2\times2\) & Uniform & 250 & 0.0306 & 1.0000 & 0.0164 & 1.0000 \\
\(2\times2\) & Uniform & 500 & 0.0408 & 1.0000 & 0.0298 & 1.0000 \\
\(2\times2\) & Uniform & 1000 & 0.0445 & 1.0000 & 0.0386 & 1.0000 \\
\addlinespace
\(2\times2\) & Same skew & 2 & 0.0000 & 0.0000 & 0.0000 & 0.0000 \\
\(2\times2\) & Same skew & 5 & 0.1110 & 0.7282 & 0.0000 & 0.2304 \\
\(2\times2\) & Same skew & 10 & 0.0357 & 0.9585 & 0.0092 & 0.6511 \\
\(2\times2\) & Same skew & 20 & 0.0232 & 0.9980 & 0.0135 & 0.9550 \\
\(2\times2\) & Same skew & 50 & 0.0086 & 1.0000 & 0.0037 & 1.0000 \\
\(2\times2\) & Same skew & 100 & 0.0107 & 1.0000 & 0.0047 & 1.0000 \\
\(2\times2\) & Same skew & 250 & 0.0260 & 1.0000 & 0.0135 & 1.0000 \\
\(2\times2\) & Same skew & 500 & 0.0379 & 1.0000 & 0.0268 & 1.0000 \\
\(2\times2\) & Same skew & 1000 & 0.0437 & 1.0000 & 0.0372 & 1.0000 \\
\addlinespace
\(2\times2\) & Different skew & 2 & 0.0000 & 0.0000 & 0.0000 & 0.0000 \\
\(2\times2\) & Different skew & 5 & 0.1008 & 0.8454 & 0.0000 & 0.3423 \\
\(2\times2\) & Different skew & 10 & 0.0432 & 0.9791 & 0.0120 & 0.7575 \\
\(2\times2\) & Different skew & 20 & 0.0244 & 0.9983 & 0.0151 & 0.9749 \\
\(2\times2\) & Different skew & 50 & 0.0080 & 1.0000 & 0.0039 & 0.9999 \\
\(2\times2\) & Different skew & 100 & 0.0121 & 1.0000 & 0.0054 & 1.0000 \\
\(2\times2\) & Different skew & 250 & 0.0290 & 1.0000 & 0.0158 & 1.0000 \\
\(2\times2\) & Different skew & 500 & 0.0387 & 1.0000 & 0.0283 & 1.0000 \\
\(2\times2\) & Different skew & 1000 & 0.0453 & 1.0000 & 0.0380 & 1.0000 \\
\addlinespace
\(3\times3\) & Uniform & 2 & 0.0000 & 0.0000 & 0.0000 & 0.0000 \\
\(3\times3\) & Uniform & 5 & 0.1931 & 0.9994 & 0.0000 & 0.9534 \\
\(3\times3\) & Uniform & 10 & 0.1058 & 1.0000 & 0.0285 & 0.9981 \\
\(3\times3\) & Uniform & 20 & 0.0730 & 1.0000 & 0.0343 & 1.0000 \\
\(3\times3\) & Uniform & 50 & 0.0288 & 1.0000 & 0.0173 & 1.0000 \\
\(3\times3\) & Uniform & 100 & 0.0167 & 1.0000 & 0.0106 & 1.0000 \\
\(3\times3\) & Uniform & 250 & 0.0269 & 1.0000 & 0.0178 & 1.0000 \\
\(3\times3\) & Uniform & 500 & 0.0355 & 1.0000 & 0.0278 & 1.0000 \\
\(3\times3\) & Uniform & 1000 & 0.0445 & 1.0000 & 0.0393 & 1.0000 \\
\addlinespace
\(3\times3\) & Same skew & 2 & 0.0000 & 0.0000 & 0.0000 & 0.0000 \\
\(3\times3\) & Same skew & 5 & 0.1807 & 0.7217 & 0.0000 & 0.2209 \\
\(3\times3\) & Same skew & 10 & 0.0602 & 0.9599 & 0.0176 & 0.6627 \\
\(3\times3\) & Same skew & 20 & 0.0333 & 0.9996 & 0.0207 & 0.9565 \\
\(3\times3\) & Same skew & 50 & 0.0135 & 1.0000 & 0.0076 & 1.0000 \\
\(3\times3\) & Same skew & 100 & 0.0139 & 1.0000 & 0.0068 & 1.0000 \\
\(3\times3\) & Same skew & 250 & 0.0273 & 1.0000 & 0.0149 & 1.0000 \\
\(3\times3\) & Same skew & 500 & 0.0348 & 1.0000 & 0.0235 & 1.0000 \\
\(3\times3\) & Same skew & 1000 & 0.0391 & 1.0000 & 0.0314 & 1.0000 \\
\addlinespace
\(3\times3\) & Different skew & 2 & 0.0000 & 0.0000 & 0.0000 & 0.0000 \\
\(3\times3\) & Different skew & 5 & 0.2089 & 0.8383 & 0.0000 & 0.3296 \\
\(3\times3\) & Different skew & 10 & 0.0886 & 0.9886 & 0.0260 & 0.7792 \\
\(3\times3\) & Different skew & 20 & 0.0466 & 0.9998 & 0.0251 & 0.9745 \\
\(3\times3\) & Different skew & 50 & 0.0188 & 1.0000 & 0.0098 & 1.0000 \\
\(3\times3\) & Different skew & 100 & 0.0171 & 1.0000 & 0.0083 & 1.0000 \\
\(3\times3\) & Different skew & 250 & 0.0268 & 1.0000 & 0.0147 & 1.0000 \\
\(3\times3\) & Different skew & 500 & 0.0331 & 1.0000 & 0.0232 & 1.0000 \\
\(3\times3\) & Different skew & 1000 & 0.0411 & 1.0000 & 0.0348 & 1.0000 \\
\addlinespace
\(5\times5\) & Uniform & 2 & 0.0000 & 0.0000 & 0.0000 & 0.0000 \\
\(5\times5\) & Uniform & 5 & 0.2287 & 1.0000 & 0.0000 & 0.9909 \\
\(5\times5\) & Uniform & 10 & 0.1061 & 1.0000 & 0.0665 & 0.9997 \\
\(5\times5\) & Uniform & 20 & 0.0768 & 1.0000 & 0.0528 & 1.0000 \\
\(5\times5\) & Uniform & 50 & 0.0643 & 1.0000 & 0.0467 & 1.0000 \\
\(5\times5\) & Uniform & 100 & 0.0377 & 1.0000 & 0.0289 & 1.0000 \\
\(5\times5\) & Uniform & 250 & 0.0232 & 1.0000 & 0.0190 & 1.0000 \\
\(5\times5\) & Uniform & 500 & 0.0276 & 1.0000 & 0.0240 & 1.0000 \\
\(5\times5\) & Uniform & 1000 & 0.0326 & 1.0000 & 0.0296 & 1.0000 \\
\addlinespace
\(5\times5\) & Same skew & 2 & 0.0000 & 0.0000 & 0.0000 & 0.0000 \\
\(5\times5\) & Same skew & 5 & 0.2175 & 0.7137 & 0.0000 & 0.2149 \\
\(5\times5\) & Same skew & 10 & 0.0825 & 0.9599 & 0.0363 & 0.6756 \\
\(5\times5\) & Same skew & 20 & 0.0651 & 0.9996 & 0.0454 & 0.9553 \\
\(5\times5\) & Same skew & 50 & 0.0290 & 1.0000 & 0.0178 & 1.0000 \\
\(5\times5\) & Same skew & 100 & 0.0112 & 1.0000 & 0.0066 & 1.0000 \\
\(5\times5\) & Same skew & 250 & 0.0152 & 1.0000 & 0.0094 & 1.0000 \\
\(5\times5\) & Same skew & 500 & 0.0369 & 1.0000 & 0.0260 & 1.0000 \\
\(5\times5\) & Same skew & 1000 & 0.0403 & 1.0000 & 0.0329 & 1.0000 \\
\addlinespace
\(5\times5\) & Different skew & 2 & 0.0000 & 0.0000 & 0.0000 & 0.0000 \\
\(5\times5\) & Different skew & 5 & 0.2690 & 0.8356 & 0.0000 & 0.3270 \\
\(5\times5\) & Different skew & 10 & 0.1514 & 0.9889 & 0.0693 & 0.7817 \\
\(5\times5\) & Different skew & 20 & 0.1259 & 1.0000 & 0.0839 & 0.9785 \\
\(5\times5\) & Different skew & 50 & 0.0623 & 1.0000 & 0.0348 & 1.0000 \\
\(5\times5\) & Different skew & 100 & 0.0250 & 1.0000 & 0.0136 & 1.0000 \\
\(5\times5\) & Different skew & 250 & 0.0254 & 1.0000 & 0.0162 & 1.0000 \\
\(5\times5\) & Different skew & 500 & 0.0360 & 1.0000 & 0.0268 & 1.0000 \\
\(5\times5\) & Different skew & 1000 & 0.0386 & 1.0000 & 0.0319 & 1.0000 \\
\addlinespace
\(8\times8\) & Uniform & 2 & 0.0000 & 0.0000 & 0.0000 & 0.0000 \\
\(8\times8\) & Uniform & 5 & 0.2692 & 0.9983 & 0.0000 & 0.9163 \\
\(8\times8\) & Uniform & 10 & 0.1376 & 1.0000 & 0.0762 & 0.9996 \\
\(8\times8\) & Uniform & 20 & 0.0687 & 1.0000 & 0.0570 & 1.0000 \\
\(8\times8\) & Uniform & 50 & 0.0703 & 1.0000 & 0.0592 & 1.0000 \\
\(8\times8\) & Uniform & 100 & 0.0712 & 1.0000 & 0.0618 & 1.0000 \\
\(8\times8\) & Uniform & 250 & 0.0317 & 1.0000 & 0.0286 & 1.0000 \\
\(8\times8\) & Uniform & 500 & 0.0212 & 1.0000 & 0.0198 & 1.0000 \\
\(8\times8\) & Uniform & 1000 & 0.0228 & 1.0000 & 0.0211 & 1.0000 \\
\addlinespace
\(8\times8\) & Same skew & 2 & 0.0000 & 0.0000 & 0.0000 & 0.0000 \\
\(8\times8\) & Same skew & 5 & 0.2266 & 0.6984 & 0.0000 & 0.2082 \\
\(8\times8\) & Same skew & 10 & 0.0873 & 0.9587 & 0.0456 & 0.6796 \\
\(8\times8\) & Same skew & 20 & 0.0873 & 0.9992 & 0.0692 & 0.9541 \\
\(8\times8\) & Same skew & 50 & 0.0537 & 1.0000 & 0.0367 & 1.0000 \\
\(8\times8\) & Same skew & 100 & 0.0245 & 1.0000 & 0.0155 & 1.0000 \\
\(8\times8\) & Same skew & 250 & 0.0106 & 1.0000 & 0.0071 & 1.0000 \\
\(8\times8\) & Same skew & 500 & 0.0162 & 1.0000 & 0.0129 & 1.0000 \\
\(8\times8\) & Same skew & 1000 & 0.0341 & 1.0000 & 0.0289 & 1.0000 \\
\addlinespace
\(8\times8\) & Different skew & 2 & 0.0000 & 0.0000 & 0.0000 & 0.0000 \\
\(8\times8\) & Different skew & 5 & 0.2986 & 0.8319 & 0.0000 & 0.3157 \\
\(8\times8\) & Different skew & 10 & 0.1853 & 0.9897 & 0.1035 & 0.7912 \\
\(8\times8\) & Different skew & 20 & 0.1803 & 1.0000 & 0.1432 & 0.9759 \\
\(8\times8\) & Different skew & 50 & 0.1391 & 1.0000 & 0.1080 & 1.0000 \\
\(8\times8\) & Different skew & 100 & 0.0921 & 1.0000 & 0.0699 & 1.0000 \\
\(8\times8\) & Different skew & 250 & 0.0583 & 1.0000 & 0.0449 & 1.0000 \\
\(8\times8\) & Different skew & 500 & 0.0377 & 1.0000 & 0.0300 & 1.0000 \\
\(8\times8\) & Different skew & 1000 & 0.0356 & 1.0000 & 0.0301 & 1.0000 \\
\addlinespace
\bottomrule\end{longtable}
\endgroup


\section{Selected additional square-table curves}

The following curves add the intermediate \(3\times3\) and \(5\times5\)
shapes to the \(2\times2\) and \(8\times8\) comparisons in
Chapter~\ref{chap:results}. The first uses uniform margins; the second
compares different margins. Upper panels use a 0--1 rejection scale and
lower panels enlarge 0--0.15. Hollow markers indicate validity below 90\%.

\mainlandscapefigure{3x3}{\(3\times3\)}{uniform}{both \(P\) and \(Q\) have uniform row and column margins \((1/3,1/3,1/3)\)}{uniform}{At \(n=10\), Wald's zero-difference rate is above 0.05 and Expanded Welch brings it below the line; both methods are almost always valid. At \(n=100\), however, both start well below 0.05 and the correction lowers rejection further. Detection grows strongly at \(n=1{,}000\), where Wald remains slightly higher. Whether the adjustment helps calibration changes with sample size, not just table shape.}
\mainlandscapefigure{5x5}{\(5\times5\)}{different_skew}{the margins are \((0.7,0.075,0.075,0.075,0.075)\) for \(P\) and \((0.8,0.05,0.05,0.05,0.05)\) for \(Q\), each used for both rows and columns}{different skew}{The \(n=10\) null is liberal for both methods, though less so for Expanded Welch; its hollow markers also show reduced validity. At \(n=100\), rejection stays low and falls across this short alternative range. At \(n=1{,}000\), detection rises and Wald remains slightly higher. Expanded Welch lowers rejection at every size, helping only where the null rate is too high.}

\FloatBarrier

\section{Other controlled comparisons}

The following comparisons change association location, sample allocation,
table shape or population construction. The final figure starts at exact
independence, where the first-order MI variance is zero.

Figure~\ref{fig:patterns-focus} compares the location of association in
\(5\times5\) tables. Columns place
association in the first \(2\times2\) block, last (rare) \(2\times2\)
block, or a centred score outer product spread across all cells. For every
column, \(P\)'s row and column margins are
\((0.7,0.075,0.075,0.075,0.075)\), and \(Q\)'s are
\((0.8,0.05,0.05,0.05,0.05)\). Here \(I(P)=0.0001\),
\(I(Q)=I(P)+|\Delta|\), \(n_P=n_Q=100\), and
\(|\Delta|=0,0.0001,0.00025,0.0005,0.001,0.002\) nats.

\begin{figure}[H]
\centering
\includegraphics{patterns_5x5_different_skew.pdf}
\caption[Changed-cell comparison, \(5\times5\) tables]{Rejection rates for
\(5\times5\) tables with association in different cells.}
\label{fig:patterns-focus}
\end{figure}

All three association patterns give below-target null rejection in this
comparison. For first, rare and spread patterns, respectively, the
zero-difference Wald rates are \(\PatternFirstWaldRate\),
\(\PatternRareWaldRate\) and \(\PatternSpreadWaldRate\); Expanded Welch gives
\(\PatternFirstExpandedRate\), \(\PatternRareExpandedRate\) and
\(\PatternSpreadExpandedRate\). All displayed calculations are valid. The
alternative curves remain nearly flat over this short MI range, so neither
method detects these differences well. The low null rejection and limited
detection persist across the three association patterns.

\clearpage

Figure~\ref{fig:imbalance-q} compares unequal samples when \(Q\) has the
higher MI. Columns have
\((n_P,n_Q)=(50,250)\) and \((250,50)\). For both columns, \(P\) has
\(3\times3\) row and column margins \((0.7,0.15,0.15)\) and \(Q\) has
\((0.8,0.1,0.1)\); the first-cell change block is used.
\(I(P)=0.02\) and \(I(Q)=0.02+|\Delta|\), with
\(|\Delta|=0,0.001,0.002,0.005,0.01,0.02\) nats.

\begin{figure}[H]
\centering
\includegraphics{imbalance_3x3_q_higher.pdf}
\caption[Unequal samples, \(Q\) higher]{Rejection rates under two sample
allocations when \(Q\) has the higher MI.}
\label{fig:imbalance-q}
\end{figure}

Allocating more observations to \(Q\), the population with higher MI,
gives stronger detection in this comparison. At zero difference, Wald is near or above 0.05 while
Expanded Welch is below it in both allocations. When \(Q\) has more observations
\((n_P,n_Q)=(50,250)\), rejection rises with \(I(Q)-I(P)\); at 0.02 nats it
reaches \(\ImbalanceQHigherWaldRate\) for Wald and
\(\ImbalanceQHigherExpandedRate\) for Expanded Welch. With \(P\) larger, the
curves rise much less. These points are valid, but the alternative gap is not
a comparison at matched false-positive rates: the null rates differ even
though both tests use the nominal 5\% threshold.

\clearpage

Figure~\ref{fig:imbalance-p} reverses the MI direction. Columns have
\((n_P,n_Q)=(50,250)\) and \((250,50)\). \(P\)'s \(3\times3\) row and
column margins are \((0.7,0.15,0.15)\); \(Q\)'s are \((0.8,0.1,0.1)\).
Both use the first-cell change block. \(I(Q)=0.02\) and
\(I(P)=0.02+|\Delta|\), with
\(|\Delta|=0,0.001,0.002,0.005,0.01,0.02\) nats.

\begin{figure}[H]
\centering
\includegraphics{imbalance_3x3_p_higher.pdf}
\caption[Unequal samples, \(P\) higher]{Rejection rates under two sample
allocations when \(P\) has the higher MI.}
\label{fig:imbalance-p}
\end{figure}

Reversing the MI direction also reverses which sample allocation gives
stronger detection. The zero-difference rates are the same as in
Figure~\ref{fig:imbalance-q}, because each allocation uses the same pair of
populations in both figures at that point; \(P\) and \(Q\) still have different
margins. Reversing the MI direction changes what happens next: when \(P\) has
the larger sample, Wald rises well above its null rate, whereas it stays
nearly flat when \(Q\) has the larger sample. At \((n_P,n_Q)=(50,250)\) and
0.02 nats, Wald rejects
\(\ImbalancePHigherWaldRate\) and Expanded Welch
\(\ImbalancePHigherExpandedRate\); both calculations are valid. Direction and
allocation therefore matter even at the same absolute MI difference.

\clearpage

Figure~\ref{fig:rectangular-focus} compares rectangular tables.
Columns are \(2\times3\) and \(3\times5\).
In each dimension \(k\), \(P\)'s corresponding margin is \(m_k(0.7)\)
and \(Q\)'s is \(m_k(0.8)\), for both rows and columns. Thus the first
column uses \(P\) rows \((0.7,0.3)\), columns
\((0.7,0.15,0.15)\), and \(Q\) rows \((0.8,0.2)\), columns
\((0.8,0.1,0.1)\); the second applies the same formula to 3 and 5
categories. Both use first-cell changes, \(n_P=n_Q=100\),
\(I(P)=0.02\), and \(I(Q)=0.02+|\Delta|\) for
\(|\Delta|=0,0.001,0.002,0.005,0.01,0.02\) nats.

\begin{figure}[H]
\centering
\includegraphics{rectangular_different_skew.pdf}
\caption[Rectangular tables]{Rejection rates for \(2\times3\) and
\(3\times5\) tables with different margins.}
\label{fig:rectangular-focus}
\end{figure}

Both rectangular panels start below 0.05, with
Expanded Welch lower. Rejection grows only modestly by 0.02 nats and remains
below 0.05 for both methods throughout the displayed range; validity is high.
Both methods therefore detect these differences poorly. Expanded Welch
also lowers an already-conservative null rate, as in the square tables.

\clearpage

Figure~\ref{fig:construction-focus} compares additive and ordinal log-linear
constructions for \(5\times5\) tables. The left column uses the first-cell additive block;
the right fits an ordinal interaction while preserving the margins. In both,
\(P\)'s row and column margins are
\((0.7,0.075,0.075,0.075,0.075)\) and \(Q\)'s are
\((0.8,0.05,0.05,0.05,0.05)\). Both have
\(I(P)=0.02\), \(I(Q)=0.02+|\Delta|\), \(n_P=n_Q=100\), and
\(|\Delta|=0,0.001,0.002,0.005,0.01,0.02\) nats.

\begin{figure}[H]
\centering
\includegraphics{construction_5x5_different_skew.pdf}
\caption[Additive versus log-linear construction]{Rejection rates under
additive and ordinal log-linear constructions for \(5\times5\) tables.}
\label{fig:construction-focus}
\end{figure}

Both constructions give conservative null rejection despite producing
different cell probabilities. At zero difference, the additive and log-linear Wald
rates are \(\ConstructionAdditiveWaldRate\) and
\(\ConstructionLoglinearWaldRate\); Expanded Welch gives
\(\ConstructionAdditiveExpandedRate\) and
\(\ConstructionLoglinearExpandedRate\). All four are valid and below 0.05.
Rejection stays low or falls slightly as the displayed difference grows.
Matching margins and MI therefore does not ensure identical rejection
rates.

\clearpage

Figure~\ref{fig:independence-focus} starts at exact independence in
\(2\times2\) tables. Columns are
uniform margins \((0.5,0.5)\) for both populations, same skew
\((0.8,0.2)\) for both, and different skew with \(P=(0.7,0.3)\) and
\(Q=(0.8,0.2)\); each vector is used for both rows and columns. The
first-cell block is used, \(I(P)=0\), \(I(Q)=|\Delta|\),
\(n_P=n_Q=100\), and
\(|\Delta|=0,0.001,0.002,0.005,0.01,0.02\) nats.

\begin{figure}[H]
\centering
\includegraphics{independence_2x2.pdf}
\caption[Exact-independence boundary]{Rejection rates for \(2\times2\)
tables starting at exact independence.}
\label{fig:independence-focus}
\end{figure}

At the uniform zero-difference point, Wald rejects
0.00025 and Expanded Welch 0.00010 of samples; validity is high. Rejection grows little over
most of the displayed alternative range, and Expanded Welch remains lower.
At exact independence, the population first-order MI variance is zero, so the approximation underlying
both curves changes. Section~\ref{sec:independence-boundary} explains the
different reference used by a one-table independence test.

\FloatBarrier
